% 数学ノート集・共通テンプレート v2（積分論2026年度第1部の版面を参考）
% A5判。章内で定義・定理・例・注意等の番号を共通化。
% コンパイル： platex -interaction=nonstopmode math_note_template_v2.tex (2回)
%              dvipdfmx math_note_template_v2.dvi
\documentclass[a5paper,10pt,onecolumn]{jsarticle}
\usepackage[dvipdfmx,top=22mm,bottom=20mm,inner=19mm,outer=18mm,headheight=13pt,headsep=6mm]{geometry}
\usepackage{amsmath,amssymb,amsthm,mathtools}
\usepackage{enumitem,fancyhdr,titlesec}
\usepackage[dvipdfmx,bookmarks=true,colorlinks=false]{hyperref}
\usepackage{pxjahyper}

% ----- 差し替える文書情報 -----
\newcommand{\NoteTitle}{複合Poissonリスクモデルと破産確率}
\newcommand{\NoteSeries}{数学ノート集}

% ----- 文字・組版 -----
% 参考PDFの質素な明朝＋Computer Modernの印象を引き継ぎながら，
% A5判で窮屈にならないよう本文の行間を少し確保する．
\linespread{1.13}
\setlength{\parindent}{1zw}
\setlength{\parskip}{0pt}
\setlength{\abovedisplayskip}{6pt plus 2pt minus 2pt}
\setlength{\belowdisplayskip}{6pt plus 2pt minus 2pt}
\setlength{\abovedisplayshortskip}{3pt plus 2pt minus 1pt}
\setlength{\belowdisplayshortskip}{4pt plus 2pt minus 1pt}
\setlist[enumerate]{topsep=3pt,itemsep=1pt,parsep=0pt,leftmargin=2.4zw}
\titleformat{\section}{\normalfont\normalsize\bfseries}{\thesection}{0.8zw}{}
\titlespacing*{\section}{0pt}{2.0ex plus .4ex}{.9ex}
\titleformat{\subsection}{\normalfont\normalsize\bfseries}{\thesubsection}{0.7zw}{}
\titlespacing*{\subsection}{0pt}{1.5ex plus .3ex}{.6ex}
\raggedbottom

% ----- 柱とページ番号：ページ上部、罫線なし -----
\fancypagestyle{mathnote}{%
  \fancyhf{}%
  \fancyhead[L]{\footnotesize \NoteSeries\ （\NoteTitle）}%
  \fancyhead[R]{\footnotesize \thepage}%
  \renewcommand{\headrulewidth}{0pt}%
  \renewcommand{\footrulewidth}{0pt}%
}
\pagestyle{mathnote}

% ----- 節内通し番号の定義・定理群 -----
\newtheoremstyle{mathnotetheorem}
 {4.5pt}{4.5pt}{\normalfont}{}%
 {\bfseries}{}{0.65zw}%
 {\thmname{#1}\thmnumber{ #2}\thmnote{（#3）}}
\theoremstyle{mathnotetheorem}
\newtheorem{theorem}{定理}[section]
\newtheorem{definition}[theorem]{定義}
\newtheorem{lemma}[theorem]{補題}
\newtheorem{proposition}[theorem]{命題}
\newtheorem{corollary}[theorem]{系}
\newtheorem{example}[theorem]{例}
\newtheorem{remark}[theorem]{注意}
\renewcommand{\proofname}{証明}
\renewcommand{\qedsymbol}{$\square$}
\makeatletter
\renewenvironment{proof}[1][\proofname]{%
  \par\pushQED{\qed}\normalfont\topsep4\p@\@plus2\p@\relax
  \trivlist\item[\hskip\labelsep\bfseries #1]\ignorespaces
}{\popQED\endtrivlist\@endpefalse}
\makeatother

% ----- 使うノートで適宜拡張する記号 -----
\newcommand{\R}{\mathbb{R}}
\newcommand{\norm}[1]{\lVert #1\rVert}
\newcommand{\abs}[1]{\lvert #1\rvert}

\begin{document}
\section*{到達目標と前提}
複合Poisson請求モデルを定め，
安全割増係数の意味を確かめた後，
初期準備金から究極破産確率を計算できるようになる．
指数請求額の場合には上界だけでなく厳密値を用い，
指定した破産確率以下にする準備金額を設計する．
前提はPoisson分布・条件付き期待値・初等的な確率過程である．


\section{Poisson過程と複合Poisson過程}

本稿では保険会社の準備金を連続時間の確率過程として
モデル化する．
確率論の条件付き期待値，独立性，大数の強法則を前提とする．

\begin{definition}
$\lambda>0$ とし，
$W_1,W_2,\dots$ を独立同分布で
$P(W_i>t)=e^{-\lambda t}$（$t\ge0$）となる
指数分布確率変数とする．
\[
 T_n=\sum_{j=1}^n W_j,\qquad T_0=0,\qquad
 N_t=\max\{n\ge0:T_n\le t\}
\]
と置く．$N_t$ を強度 $\lambda$ の
\emph{Poisson過程}という．
\end{definition}

\begin{lemma}[Poisson分布]\label{lem:poisson}
任意の $t\ge0$，$n\ge0$ について
\[
 P(N_t=n)=e^{-\lambda t}\frac{(\lambda t)^n}{n!}
\]
が成り立つ．
さらに互いに交わらない時間区間の増分は独立であり，
長さ $h$ の区間の増分は平均 $\lambda h$ の
Poisson分布に従う．
\end{lemma}
\begin{proof}
$n\ge1$ の場合，
$T_n$ は $n$ 個の指数分布の和だから，
その密度は畳み込みを用いる帰納法により
\[
 f_{T_n}(s)
 =\frac{\lambda^n s^{n-1}e^{-\lambda s}}{(n-1)!},
 \qquad s>0
\]
となる．
実際 $n=1$ は指数分布の密度であり，
$n$ で成立すれば
\[
 f_{T_{n+1}}(t)
 =\int_0^t f_{T_n}(s)\lambda e^{-\lambda(t-s)}\,ds
 =\frac{\lambda^{n+1}t^ne^{-\lambda t}}{n!}
\]
となる．
従って
\begin{align*}
 P(N_t=n)
 &=P(T_n\le t<T_{n+1})\\
 &=\int_0^t f_{T_n}(s)P(W_{n+1}>t-s)\,ds\\
 &=\int_0^t
 \frac{\lambda^ns^{n-1}e^{-\lambda s}}{(n-1)!}
 e^{-\lambda(t-s)}\,ds\\
 &=e^{-\lambda t}\frac{(\lambda t)^n}{n!}.
\end{align*}
$n=0$ では $P(W_1>t)=e^{-\lambda t}$ である．

増分の独立性を示す．
指数分布は記憶を持たず，
\[
 P(W>s+t\mid W>s)
 =\frac{e^{-\lambda(s+t)}}{e^{-\lambda s}}
 =e^{-\lambda t}
\]
である．
従って任意の時刻 $t$ の直後の待ち時間は，
過去の待ち時間と到着回数を条件としても
強度 $\lambda$ の指数分布であり，
その後に続く待ち時間も元の独立な指数分布列と同じ法則を持つ．
これを到着時刻ごとに反復すると，
$t$ 以後の計数過程は過去から独立で，
時刻零からの過程と同じ有限次元分布を持つ．
従って定常独立増分性が成立し，
増分の分布は既に示したPoisson分布の公式による．
\end{proof}

\begin{definition}
$X_1,X_2,\dots$ を非負の独立同分布確率変数とし，
$(N_t)$ と独立であるとする．
\[
 S_t=\sum_{j=1}^{N_t}X_j
\]
を\emph{複合Poisson過程}という．
$X_j$ は一回の保険金請求額を表す．
\end{definition}

\begin{proposition}[複合Poisson過程の期待値と分散]
\label{prop:compound-moments}
$E[X_1]<\infty$ なら
\[
 E[S_t]=\lambda tE[X_1].
\]
$E[X_1^2]<\infty$ なら
\[
 \operatorname{Var}(S_t)=\lambda tE[X_1^2].
\]
\end{proposition}
\begin{proof}
$N_t=n$ のもとでは
$S_t=X_1+\cdots+X_n$ なので，
\[
 E[S_t\mid N_t]=N_tE[X_1].
\]
期待値を取れば
$E[S_t]=E[N_t]E[X_1]=\lambda tE[X_1]$ となる．
二乗モーメントが有限なら
\[
 \operatorname{Var}(S_t\mid N_t)
 =N_t\operatorname{Var}(X_1).
\]
全分散公式から
\begin{align*}
 \operatorname{Var}(S_t)
 &=E[\operatorname{Var}(S_t\mid N_t)]
 +\operatorname{Var}(E[S_t\mid N_t])\\
 &=\lambda t\operatorname{Var}(X_1)
 +\lambda t(E[X_1])^2\\
 &=\lambda tE[X_1^2].
\end{align*}
Poisson分布の平均と分散が $\lambda t$ であることは，
母関数 $\exp(\lambda t(z-1))$ を微分して確認できる．
\end{proof}

\section{準備金過程と破産時刻}

\begin{definition}
初期準備金 $u\ge0$，単位時間あたりの
保険料収入 $c>0$ に対して
\[
 U_t=u+ct-S_t
\]
を\emph{Cram\'er--Lundberg準備金過程}という．
\emph{破産時刻}を
\[
 \tau_u=\inf\{t\ge0:U_t<0\}
\]
と定義する．
\emph{究極破産確率}を
$\psi(u)=P(\tau_u<\infty)$ と書く．
\end{definition}

\begin{definition}
$c>\lambda E[X_1]$ を
\emph{正の純利益条件}という．
$c=(1+\theta)\lambda E[X_1]$ と書くとき
$\theta>0$ を安全割増係数という．
\end{definition}

\begin{proposition}[純利益条件の意味]\label{prop:net-profit}
$E[X_1]<\infty$ とすると
\[
 \frac{S_t}{t}\longrightarrow\lambda E[X_1]
 \qquad\text{ほとんど確実に}
\]
が成り立つ．
従って $c<\lambda E[X_1]$ なら，
任意の有限初期準備金 $u$ について
$\psi(u)=1$ である．
\end{proposition}
\begin{proof}
独立同分布の指数待ち時間の強法則により
$T_n/n\to1/\lambda$ がほとんど確実に成立する．
$T_{N_t}\le t<T_{N_t+1}$ と
$N_t\to\infty$ を用いると
$N_t/t\to\lambda$ がほとんど確実に成立する．
また請求額列の強法則から
\[
 \frac1{N_t}\sum_{j=1}^{N_t}X_j
 \longrightarrow E[X_1]
\]
がほとんど確実に成立する．
両式を掛けて $S_t/t\to\lambda E[X_1]$ を得る．
$c<\lambda E[X_1]$ なら
\[
 \frac{U_t}{t}=\frac ut+c-\frac{S_t}{t}
 \longrightarrow c-\lambda E[X_1]<0
\]
がほとんど確実に成立する．
従って十分大きい $t$ で $U_t<0$ となり，
破産時刻は有限である．
\end{proof}

\section{調整係数と指数マルチンゲール}

\begin{definition}
請求額の積率母関数を
\[
 M_X(r)=E[e^{rX_1}]
\]
とする．
正の実数 $R>0$ が
\[
 cR=\lambda(M_X(R)-1)
 \tag{L}
\]
を満たすとき，$R$ を\emph{Lundberg調整係数}という．
\end{definition}

\begin{lemma}[調整係数の存在と一意性]\label{lem:adjustment}
$0<E[X_1]<\infty$ とし，
$c>\lambda E[X_1]$ とする．
$M_X(r)$ がある区間 $[0,r_1]$ で有限であり，
\[
 \lambda(M_X(r_1)-1)>cr_1
\]
を満たすとする．
このとき $(0,r_1)$ にただ一つの調整係数 $R$ が存在する．
\end{lemma}
\begin{proof}
\[
 h(r)=\lambda(M_X(r)-1)-cr
\]
と置く．$h(0)=0$ であり，
右微分は $h'(0+)=\lambda E[X_1]-c<0$ である．
また $h(r_1)>0$ なので連続性と中間値の定理から
$h(R)=0$ となる $R\in(0,r_1)$ が存在する．
一意性を示す．
$P(X_1>0)>0$ だから，
$e^{rX_1}$ の $r$ に関する厳密凸性と期待値の線形性により
$M_X(r)$ は $[0,r_1]$ 上で厳密凸である．
従って $h$ も厳密凸である．
$h(0)=0$ と $h'(0+)<0$ により，
$h$ は最初負となり，正の方向で零を横切る回数は高々一回である．
実際二つの正の零点 $0<R_1<R_2$ があれば，
$R_1$ は $0$ と $R_2$ の真の凸結合であるため
厳密凸性から
$h(R_1)<(1-R_1/R_2)h(0)+(R_1/R_2)h(R_2)=0$
となり矛盾する．
\end{proof}

\begin{lemma}[複合Poisson過程の積率母関数]
\label{lem:compound-mgf}
$M_X(r)<\infty$ なら
\[
 E[e^{rS_t}]
 =\exp\{\lambda t(M_X(r)-1)\}.
\]
\end{lemma}
\begin{proof}
$N_t=n$ のもとで $X_j$ は独立なので
$E[e^{rS_t}\mid N_t=n]=M_X(r)^n$ である．
Poisson分布に関する期待値を取ると
\begin{align*}
 E[e^{rS_t}]
 &=\sum_{n=0}^\infty
 e^{-\lambda t}\frac{(\lambda t)^n}{n!}M_X(r)^n\\
 &=e^{-\lambda t}
   \exp(\lambda tM_X(r))
 =\exp\{\lambda t(M_X(r)-1)\}.
\end{align*}
\end{proof}

\begin{theorem}[指数マルチンゲール]\label{thm:exp-mart}
$R>0$ が (L) を満たし，
$M_X(R)<\infty$ とする．
このとき
\[
 Z_t=e^{-RU_t}
\]
は，請求過程の自然なフィルトレーションに関して
非負マルチンゲールであり，
$E[Z_t]=e^{-Ru}$ が全ての $t\ge0$ で成り立つ．
\end{theorem}
\begin{proof}
$U_t=u+ct-S_t$ と書くと
\[
 Z_t=e^{-Ru}e^{-Rct+rS_t}\big|_{r=R}.
\]
補題\ref{lem:compound-mgf}と (L) から
\[
 E[Z_t]
 =e^{-Ru}\exp\{-Rct+\lambda t(M_X(R)-1)\}
 =e^{-Ru}.
\]
次に $0\le s<t$ とする．
$S_t-S_s$ は $\mathcal F_s$ から独立で，
$S_{t-s}$ と同じ分布を持つので
\begin{align*}
 E[Z_t\mid\mathcal F_s]
 &=Z_sE[e^{-Rc(t-s)+R(S_t-S_s)}\mid\mathcal F_s]\\
 &=Z_s\exp\{-Rc(t-s)
       +\lambda(t-s)(M_X(R)-1)\}\\
 &=Z_s.
\end{align*}
従ってマルチンゲールである．
\end{proof}

\section{Lundbergの破産確率上界}

\begin{lemma}[非負マルチンゲールの有界停止]
\label{lem:bounded-stop}
$(Z_t)_{t\ge0}$ を右連続な非負マルチンゲール，
$\tau$ を停止時刻とする．
任意の整数 $T\ge1$ について
\[
 E[Z_{\tau\wedge T}]\le E[Z_0]
\]
が成り立つ．
\end{lemma}
\begin{proof}
$n\ge1$ に対して
$\tau_n=2^{-n}\lceil2^n(\tau\wedge T)\rceil$
と置く．
$\tau_n$ は有限個の値を取る停止時刻であり，
$0\le\tau_n\le T$ である．
離散時刻の停止定理を有限回の条件付き期待値で
示すと，$E[Z_{\tau_n}]=E[Z_0]$ となる．
実際 $2^{-n}$ 刻みの過程 $M_k=Z_{k2^{-n}}$ は
離散マルチンゲールであり，
$M_{N_n}=M_0+
 \sum_{k=1}^{N_n}1_{\{\tau_n\ge k2^{-n}\}}
 (M_k-M_{k-1})$ と書ける．
各指示関数は前時刻までの情報で決まるため，
期待値を取ると和の各項は零である．
$N_n=2^nT$ として上の等式が従う．

$\tau_n\downarrow\tau\wedge T$ なので，
$Z$ の右連続性から
$Z_{\tau_n}\to Z_{\tau\wedge T}$ が成り立つ．
$Z$ は非負なのでFatouの補題から
\[
 E[Z_{\tau\wedge T}]
 \le\liminf_{n\to\infty}E[Z_{\tau_n}]
 =E[Z_0].
\]
\end{proof}

\begin{theorem}[Lundbergの不等式]\label{thm:lundberg}
調整係数 $R>0$ が存在するなら，
任意の初期準備金 $u\ge0$ に対して
\[
 \psi(u)=P(\tau_u<\infty)\le e^{-Ru}.
\]
\end{theorem}
\begin{proof}
指数マルチンゲール $Z_t=e^{-RU_t}$ を考える．
準備金は請求到着時以外では傾き $c>0$ で増加するため，
破産するときの準備金は $U_{\tau_u}<0$ である．
従って $\{\tau_u\le T\}$ 上では
$Z_{\tau_u\wedge T}=e^{-RU_{\tau_u}}>1$ となる．
補題\ref{lem:bounded-stop}によって
\[
 P(\tau_u\le T)
 \le E[Z_{\tau_u\wedge T}]
 \le E[Z_0]=e^{-Ru}.
\]
$T\to\infty$ として
$P(\tau_u<\infty)\le e^{-Ru}$ を得る．
\end{proof}

\section{初回請求による積分方程式}

\begin{theorem}[破産確率の積分方程式]
\label{thm:ruin-integral}
$F_X$ を請求額の分布関数とする．
任意の $u\ge0$ について
\[
 \begin{aligned}
 \psi(u)=\int_0^\infty\lambda e^{-\lambda t}
 \biggl\{&
 P(X_1>u+ct)\\
 &+\int_{[0,u+ct]}
       \psi(u+ct-x)\,dF_X(x)
 \biggr\}\,dt.
 \end{aligned}
 \tag{R}
\]
\end{theorem}
\begin{proof}
最初の請求到着時刻 $T_1$ は強度 $\lambda$ の指数分布である．
その直前の準備金は $u+cT_1$ である．
請求額 $X_1>u+cT_1$ ならその瞬間に破産する．
一方 $X_1\le u+cT_1$ なら，
残余準備金は $u+cT_1-X_1\ge0$ であり，
その後の請求過程は過去から独立で元の過程と同じ法則を持つ．
従ってその条件付き究極破産確率は
$\psi(u+cT_1-X_1)$ である．
$T_1$ と $X_1$ の独立性から，
二つの場合分けの条件付き期待値を
指数分布の密度 $\lambda e^{-\lambda t}$ に関して積分すると
表示式が得られる．
\end{proof}

\begin{proposition}[減衰する有界解の一意性]
\label{prop:ruin-unique}
$c>\lambda E[X_1]$ と仮定する．
式 (R) を満たし，
$u\to\infty$ で零に収束する有界関数
$\psi:[0,\infty)\to\mathbb R$ は高々一つである．
\end{proposition}
\begin{proof}
$\psi_1,\psi_2$ が二つの解とし，
$h=\psi_1-\psi_2$ と置く．
(R) の差を取ると
\[
 h(u)=E[1_{\{X_1\le u+cW_1\}}
 h(u+cW_1-X_1)]
\]
となる．
これを独立増分と条件付き期待値で反復すると
\[
 h(u)=E[1_{\{\tau_u>T_n\}}h(U_{T_n})]
\]
を得る．
ただし $T_n=W_1+\cdots+W_n$ である．
大数の強法則から
\[
 \frac{U_{T_n}}n
 =\frac un+c\frac{T_n}n
 -\frac1n\sum_{j=1}^nX_j
 \longrightarrow\frac c\lambda-E[X_1]>0
\]
がほとんど確実に成立する．
従って $U_{T_n}\to\infty$ であり，
$\tau_u>T_n$ の指示関数が1となる経路では
$h(U_{T_n})\to0$ である．
$h$ は有界なのでLebesgueの収束定理から
\[
 h(u)=\lim_{n\to\infty}
 E[1_{\{\tau_u>T_n\}}h(U_{T_n})]=0.
\]
$u$ は任意なので二つの解は一致する．
\end{proof}

\section{指数請求額の場合の厳密解}

\begin{theorem}[指数分布の破産確率]\label{thm:exponential-ruin}
請求額 $X_i$ が率 $\beta>0$ の指数分布に従い，
$c>\lambda/\beta$ とする．
このとき
\[
 \psi(u)=\frac{\lambda}{c\beta}
 \exp\!\left[-\left(\beta-\frac\lambda c\right)u\right].
\]
\end{theorem}
\begin{proof}
指数分布の積率母関数は
$M_X(r)=\beta/(\beta-r)$（$0\le r<\beta$）である．
調整係数方程式
$cR=\lambda(M_X(R)-1)$ に代入すると
\[
 cR=\lambda\frac R{\beta-R}
\]
であり，正の解は
$R=\beta-\lambda/c>0$ となる．
$A=\lambda/(c\beta)$ と置くと
\[
 A=\frac{\beta-R}{\beta}
\]
である．
候補関数 $\phi(u)=Ae^{-Ru}$ が
式 (R) を満たすことを示す．
指数分布の密度を用いると
\begin{align*}
 &P(X>z)+\int_0^z\phi(z-x)\beta e^{-\beta x}\,dx\\
 &\quad=e^{-\beta z}
 +A\beta e^{-Rz}
   \int_0^z e^{-(\beta-R)x}\,dx\\
 &\quad=e^{-\beta z}
 +\frac{A\beta}{\beta-R}
  (e^{-Rz}-e^{-\beta z})\\
 &\quad=e^{-Rz},
\end{align*}
最後に $A\beta/(\beta-R)=1$ を用いた．
この式に $z=u+ct$ を代入して (R) の右辺に入れると
\begin{align*}
 &\int_0^\infty\lambda e^{-\lambda t}
 e^{-R(u+ct)}\,dt\\
 &\quad=e^{-Ru}\frac{\lambda}{\lambda+cR}
 =e^{-Ru}\frac{\lambda}{c\beta}
 =Ae^{-Ru}=\phi(u).
\end{align*}
従って $\phi$ は (R) を満たす．
Lundbergの不等式から真の破産確率 $\psi(u)$ も
$u\to\infty$ で零に収束する．
$\phi$ も有界で零へ収束するので，
命題\ref{prop:ruin-unique}から
$\psi=\phi$ となる．
\end{proof}

\begin{remark}
調整係数による指数上界は，
請求額の積率母関数が正の点で有限であるという仮定に依存する．
重い裾を持つ請求額分布では
正の調整係数が存在しないこともあり，
破産確率の減衰は指数型とは限らない．
\end{remark}

\section{破産確率の上限から準備金を設計する}
\begin{example}[指数請求額の具体的設計]
請求到着率を $\lambda=1$，
請求額を指数分布 $\mathrm{Exp}(\beta)$
（率 $\beta=2$，平均請求額 $1/2$），
保険料収入率を $c=1$ とする．
正の純利益条件は
$c=1>\lambda E[X_1]=1/2$ により成立する．
安全割増係数は
$\theta=c/(\lambda E[X_1])-1=1$ である．
指数請求額の場合の厳密公式から調整係数
$R=\beta-\lambda/c=1$ を得るので
\[
 \psi(u)=\frac{\lambda}{c\beta}e^{-Ru}
        =\frac12 e^{-u}\qquad (u\ge0).
\]
特に初期準備金が零なら $\psi(0)=1/2$，
$u=\log5$ とすれば $\psi(u)=1/10$ である．
破産確率を $5\%$ 以下にする条件は
\[
 \psi(u)\le\frac1{20}
 \quad\Longleftrightarrow\quad
 \frac12 e^{-u}\le\frac1{20}
 \quad\Longleftrightarrow\quad u\ge\log10.
\]
このモデルでは $\log10$ が必要十分な初期準備金である．
一方，Lundbergの上界 $\psi(u)\le e^{-Ru}=e^{-u}$
だけを使うと $u\ge\log20$ を要求するため，
上界から設計した場合には余裕のある値となる．
\end{example}
\begin{proposition}[指数請求額の一般的な設計式]
請求到着率 $\lambda>0$，指数請求額の率
$\beta>0$，保険料率 $c>\lambda/\beta$ とする．
初期準備金を用いて目標破産確率
$\varepsilon\in(0,1)$ を達成するための必要十分条件は
\[
 u\ge \max\left\{0,\,
 \frac{\log(\lambda/(c\beta\varepsilon))}
      {\beta-\lambda/c}\right\}.
\]
\end{proposition}
\begin{proof}
既に示した厳密解
$\psi(u)=(\lambda/(c\beta))e^{-Ru}$，
$R=\beta-\lambda/c>0$ を使う．
指数関数は狭義減少だから $\psi(u)\le\varepsilon$ と
$u\ge R^{-1}\log(\lambda/(c\beta\varepsilon))$
は同値である．準備金は $u\ge0$ に制限するので，
両条件の共通部分が表示式となる．
\end{proof}
\section*{到達確認問題と解答}
\noindent\textbf{問題．}
最初の例で破産確率を $1\%$ 以下にする
最低の初期準備金を求めよ．

\medskip\noindent\textbf{解答．}
$(1/2)e^{-u}\le0.01$ は
$e^{-u}\le0.02$ に同値なので $u\ge\log50$．
従って最低準備金は $\log50$ である．
\section*{仮定の意味}
以上の数値は指数分布・独立請求・定率保険料という
仮定の下の計算であり，
重い裾の請求分布や請求間の依存を許す場合には
同じ公式を適用できない．

\end{document}
